\chapter{MRSA}
\index{MRSA}

\medskip
\noindent\textit{\textbf{Under review.} This chapter is an AI-assisted educational draft; it carries no source citations and is not a substitute for professional medical advice.}

\section*{Definition and Overview}
Methicillin-resistant \textit{Staphylococcus aureus} (MRSA) is a strain of \textit{S. aureus} resistant to beta-lactam antibiotics (penicillins, cephalosporins, carbapenems) due to acquisition of the \textit{mecA} gene (or \textit{mecC}) on the staphylococcal cassette chromosome \textit{mec} (SCC\textit{mec}). It causes skin and soft tissue infections, pneumonia, bloodstream infections, and toxic shock syndrome. Classified as healthcare-associated (HA-MRSA), community-associated (CA-MRSA), or livestock-associated (LA-MRSA). CA-MRSA typically carries SCC\textit{mec} type IV/V and Panton-Valentine leukocidin (PVL), causing severe necrotising infections.

\section*{Epidemiology}
In many countries, MRSA accounts for 15--20\% of \textit{S. aureus} bacteraemias. HA-MRSA rates have declined with infection control; CA-MRSA is increasing. Indigenous people have higher CA-MRSA rates (clonal complex 93, ST93). Risk factors: hospitalisation, surgery, dialysis, indwelling devices, immunosuppression, contact sports, crowded living, injecting drug use. Global burden: ~100,000 invasive MRSA infections annually in US; rising in Asia-Pacific.

\section*{Aetiology and Pathophysiology}
\begin{itemize}
    \item \textbf{Resistance mechanism:} \textit{mecA} encodes PBP2a (penicillin-binding protein 2a) with low affinity for beta-lactams; allows cell wall synthesis despite antibiotic presence
    \item \textbf{SCC\textit{mec} types:} I--III (HA-MRSA, larger, multi-resistant); IV--V (CA-MRSA, smaller, often PVL+); XIII (mecC-MRSA, zoonotic)
    \item \textbf{Virulence factors:} PVL (neutrophil lysis, tissue necrosis), alpha-toxin, TSST-1 (toxic shock), enterotoxins, biofilm formation, protein A, clumping factors
    \item \textbf{Transmission:} direct contact, contaminated surfaces, healthcare workers' hands, fomites; nasal colonisation (~30\% population, ~1--3\% MRSA)
    \item \textbf{Pathogenesis:} colonisation $\to$ breach of skin/mucosa $\to$ invasion $\to$ toxin-mediated tissue destruction $\to$ haematogenous spread $\to$ metastatic infection (osteomyelitis, endocarditis, septic arthritis)
\end{itemize}

\section*{Clinical Features}
\begin{itemize}
    \item \textbf{Skin/soft tissue:} abscess, furuncle, carbuncle, cellulitis, necrotising fasciitis (PVL+ CA-MRSA); often "spider bite" appearance
    \item \textbf{Pneumonia:} necrotising, cavitary, haemorrhagic; post-influenza; high mortality (30--50\%)
    \item \textbf{Bacteraemia/endocarditis:} fever, rigors, metastatic foci; right-sided endocarditis in IVDU
    \item \textbf{Bone/joint:} osteomyelitis, septic arthritis (especially children)
    \item \textbf{Toxic shock:} fever, rash, hypotension, multi-organ failure (TSST-1)
    \item \textbf{Colonisation:} asymptomatic nasal, perineal, axillary carriage
\end{itemize}

\section*{Differential Diagnosis}
\begin{itemize}
    \item MSSA (methicillin-susceptible \textit{S. aureus}) infections -- identical clinical picture; distinguished by susceptibility testing
    \item \textit{Streptococcal} infections (cellulitis, necrotising fasciitis) -- usually beta-haemolytic strep; no catalase
    \item Gram-negative infections (\textit{Pseudomonas}, Enterobacteriaceae) -- different risk factors, empiric coverage differs
    \item \textit{Clostridial} myonecrosis -- gas in tissues, crepitus, rapid progression
    \item Non-infectious: vasculitis, malignancy, gout (for septic arthritis), inflammatory bowel disease (for abscess)
\end{itemize}

\section*{When to Seek Medical Care}
Seek urgent care for rapidly spreading skin infection, fever with skin lesions, signs of sepsis, or infection not improving on standard antibiotics.

\textbf{Contact your local emergency services for:}
\begin{itemize}
    \item Hypotension, altered consciousness, respiratory distress (sepsis/toxic shock)
    \item Rapidly expanding painful red skin with crepitus or necrosis (necrotising infection)
    \item Haemoptysis with cavitary pneumonia
\end{itemize}

\section*{Diagnosis and Investigations}
\begin{itemize}
    \item \textbf{Culture:} swab (nasal, wound, sputum, blood); chromogenic agar for MRSA screening; 18--24h
    \item \textbf{PCR:} rapid \textit{mecA}/\textit{mecC} + \textit{spa} typing; 1--2h (e.g., GeneXpert, BD MAX)
    \item \textbf{Antibiotic susceptibility:} cefoxitin disk diffusion (surrogate for mecA); MIC for vancomycin, daptomycin, linezolid, ceftaroline
    \item \textbf{Blood tests:} FBC, CRP, procalcitonin, creatinine, LFTs, coagulation, blood cultures $\times$2
    \item \textbf{Imaging:} CXR (pneumonia); echo (endocarditis); MRI/CT (deep abscess, osteomyelitis); ultrasound (joint effusion)
    \item \textbf{Screening:} nasal swab pre-op, ICU admission, high-risk units; decolonisation if positive
\end{itemize}

\section*{Management}
\begin{itemize}
    \item \textbf{Source control:} incision and drainage (I\&D) for abscesses -- often sufficient for small CA-MRSA skin infections; surgical debridement for necrotising infection
    \item \textbf{Non-severe skin/soft tissue (CA-MRSA):} oral clindamycin, doxycycline/minocycline, TMP-SMX, linezolid; 7--10 days; consider rifampin adjunct for biofilm
    \item \textbf{Severe/HA-MRSA (bacteraemia, pneumonia, deep infection):} IV vancomycin (target AUC 400--600 mg$\cdot$h/L) or daptomycin 8--10 mg/kg (not for pneumonia); alternative: linezolid, ceftaroline, telavancin
    \item \textbf{MRSA pneumonia:} linezolid or ceftaroline preferred (lung penetration); vancomycin if susceptible
    \item \textbf{Endocarditis:} vancomycin + gentamicin (short course) +/- rifampin; 6 weeks native valve
    \item \textbf{Decolonisation:} mupirocin 2\% nasal ointment BD $\times$5 days + chlorhexidine 4\% body wash daily $\times$5 days; household contacts simultaneously
    \item \textbf{Infection control:} contact precautions, single room, hand hygiene, environmental cleaning, antimicrobial stewardship
\end{itemize}

\section*{Complications}
\begin{itemize}
    \item Metastatic seeding: endocarditis, osteomyelitis, septic arthritis, epidural abscess, brain abscess
    \item Sepsis, septic shock, multi-organ failure
    \item Necrotising pneumonia with empyema, pneumatocele
    \item Toxic shock syndrome (TSST-1)
    \item Vancomycin treatment failure (MIC creep, heteroresistance); nephrotoxicity, thrombocytopenia
    \item Recurrence (persistent colonisation, inadequate source control)
\end{itemize}

\section*{Prognosis and Outcomes}
CA-MRSA skin infections: excellent with I\&D + appropriate oral therapy. HA-MRSA bacteraemia: 30-day mortality 20--30\%. MRSA pneumonia: mortality 30--50\%. Endocarditis: mortality 30--40\%. PVL+ strains associated with severe necrotising disease. Early source control and appropriate antibiotics improve outcomes. Recurrence risk 10--20\% within 1 year.

\section*{Prevention and Self-Care}
\begin{itemize}
    \item Hand hygiene (alcohol-based rub or soap/water)
    \item Cover wounds; don't share towels, razors, sports equipment
    \item Decolonisation for known carriers pre-surgery/ICU
    \item Antimicrobial stewardship: avoid unnecessary beta-lactams/fluoroquinolones
    \item Contact precautions in healthcare; screen high-risk admissions
    \item Environmental cleaning (high-touch surfaces, shared equipment)
    \item Vaccination (influenza, pneumococcal) reduces superinfection risk
\end{itemize}

\section*{Sources and Further Reading}
\begin{itemize}
    \item Therapeutic Guidelines (Australia). \textit{Antibiotic guidelines -- Staphylococcus aureus (including MRSA).} https://www.tg.org.au/
    \item Australian Commission on Safety and Quality in Health Care. \textit{MRSA prevention and control.} https://www.safetyandquality.gov.au/
    \item Centers for Disease Control and Prevention (CDC). \textit{Methicillin-resistant Staphylococcus aureus (MRSA) infections.} https://www.cdc.gov/mrsa/
    \item Infectious Diseases Society of America (IDSA). \textit{MRSA treatment guidelines.} https://www.idsociety.org/
    \item Healthdirect Australia. \textit{MRSA (golden staph).} https://www.healthdirect.gov.au/mrsa
    \item Turnidge JD et al. \textit{Staphylococcus aureus bacteraemia in Australian hospitals.} MJA 2019.
    \item Tong SYC et al. \textit{Staphylococcus aureus infections: epidemiology, pathophysiology, clinical manifestations, and management.} Clin Microbiol Rev 2015.
\end{itemize}

\clearpage